\subsection{局部有限族}

定义 8. 族 \((\mathbf{A}_{i})_{i\in \mathbf{I}}\)（拓扑空间 \(\mathbf{X}\) 的子集所成的族）称为局部有限的，如果对每个 \(x\in \mathbf{X}\)，存在一个邻域 \(\mathbf{V}\)（\(x\) 的邻域），使得 \(\mathbf{V}\cap \mathbf{A}_i = \emptyset\) 对除去有限多个指标 \(i\in \mathbf{I}\) 外成立。子集所成的集合 \(\mathfrak{S}\)（\(\mathbf{X}\) 的子集所成之集）称为局部有限的，如果由 \(\mathfrak{S}\) 到自身的恒等映射所定义的子集族是局部有限的。

显然，如果 \((\mathbf{A}_t)_{t \in \mathbf{I}}\) 是局部有限的子集族，并且 \(\mathbf{B}_t \subset \mathbf{A}_t\) 对每个 \(t \in \mathbf{I}\) 成立，则族 \((\mathbf{B}_t)_{t \in \mathbf{I}}\) 是局部有限的。

拓扑空间的每个有限子集族显然是局部有限的；反之一般不成立。

* 例如，在 R 中，由区间 \(]\leftarrow, 1[\) 与诸区间 \(]n, \rightarrow[\)（对每个整数 \(n\geqslant0\)）所形成的开覆盖是局部有限的；而每个区间 \(]n, \rightarrow[\) 都与该覆盖中的无穷多个集合相交。

命题 4. 拓扑空间 X 的闭子集所成的局部有限族之并在 X 中是闭集。

设 \((\mathbf{F}_i)_{i\in I}\) 是 \(X\) 的闭子集所成的一个局部有限族，并设 \(x\in X\) 不属于 \(\mathbf{F} = \bigcup_{i\in I}\mathbf{F}_i\)；则 \(x\) 有一个邻域 \(V\)，它只与那些 \(\mathbf{F}_i\) 相交，后者的指标属于某个有限子集 \(J\)（\(I\) 的子集）。对每个 \(i\in J\)，设 \(\mathbf{U}_i\) 为 \(\mathbf{F}_i\) 的余集；则 \(\mathbb{CF}\) 包含集合 \(V\cap \bigcap_{i\in J}\mathbf{U}_i\)，后者是 \(x\) 的一个邻域，因为每个 \(\mathbf{U}_i\) 都是开集且包含 \(x\)。于是，由第 2 目的命题 1，\(\mathbb{CF}\) 是开集，从而 \(\mathbf{F}\) 在 \(X\) 中是闭集。

我们注意到，X 的闭子集所成的任意族之并未必是闭集；例如，在有理直线 Q 上，集合 \(]0, 1[\) 是诸闭集

\[
\left[ \frac {1}{n}, 1 - \frac {1}{n} \right] \quad \text { for } n > 2,
\]

之并，但它不是闭集。

